% =====================================================================
% chapters/02_literature_review.tex
% Chapter 2 — Literature Review
% Project: Vision-Based Running Technique Classification from In-the-Wild
%          Internet Videos using Markerless Pose Estimation and Deep Learning
%
% Required preamble packages (declare in main.tex):
%   \usepackage{booktabs}  \usepackage{array}  \usepackage{tabularx}
%
% Citation keys are drawn from references.bib and from the verified papers in
% the project literature folder. No empirical results appear here.
% =====================================================================

\chapter{Literature Review}
\label{ch:literature-review}

This chapter reviews the research relevant to vision-based running-technique
classification. It proceeds from the biomechanics of running form, through the
sensing modalities used to measure it, to the pose-estimation, cycle-detection,
and classification methods on which an automatic pipeline can be built. It then
examines how multi-cycle video data should be evaluated, and closes by stating
the research gap that this thesis addresses.

% ---------------------------------------------------------------------
\section{Running Technique Analysis}
\label{sec:lit-running}

Running technique, or running form, is commonly characterised by a small set of
inter-related kinematic descriptors: the manner of \emph{foot strike} (rearfoot,
midfoot, or forefoot contact), trunk and pelvis \emph{posture}, \emph{stride}
parameters such as step length and cadence, and the \emph{coordination} of the
upper and lower limbs through the gait cycle
\cite{novacheck1998biomechanics, inostroza2025runningcycle}. These
descriptors are not independent; for example, an over-strided rearfoot contact
typically coincides with a more extended knee at touchdown and altered trunk
lean, so that "correct" and "incorrect" technique are better understood as
configurations of the whole kinematic chain than as single isolated events
\cite{dicharry2010kinematics}.

A recurring theme in the running-biomechanics literature is that technique is
expressed \emph{periodically}: the discriminative information lives in how the
body moves through repeated running cycles rather than in any single static
frame \cite{inostroza2025runningcycle, cutting1978gaitperiodicity}. This has two
consequences for an automatic system. First, the unit of analysis should be a
running cycle, not an arbitrary frame window. Second, the descriptors most
diagnostic of form --- foot contact, knee and ankle behaviour at touchdown, limb
swing --- are concentrated in the lower body, which motivates particular
attention to lower-limb measurement (Section~\ref{sec:lit-pose}).

% ---------------------------------------------------------------------
\section{Marker-Based and Markerless Motion Analysis}
\label{sec:lit-modalities}

Three broad modalities are used to capture running kinematics:
optical \emph{marker-based} motion capture, \emph{wearable} inertial sensors,
and \emph{markerless} video-based pose estimation. Marker-based systems track
reflective markers with calibrated multi-camera rigs and are widely regarded as
the laboratory reference for kinematic accuracy, but they require controlled
environments, specialist hardware, and marker placement, which makes them
unsuitable for analysing the large volumes of uncontrolled footage found online
\cite{panconi2024markerless}. Wearable inertial measurement units estimate
segment motion from accelerometers and gyroscopes; they are portable and usable
outdoors, but require body-worn instrumentation, calibration, and drift
correction, and are likewise absent from existing internet video
\cite{fiaz2026gaitphase, bulling2014tutorial}.

Markerless video-based analysis estimates body configuration directly from
ordinary RGB video using learned pose-estimation models. It requires no
instrumentation and can in principle be applied to any recording, which is what
makes it the only viable modality for \emph{in-the-wild} internet video. The
trade-off is reduced and less certain measurement: estimates are noisier,
subject to occlusion and viewpoint variation, and accompanied only by
model-internal confidence rather than physical calibration
\cite{lugaresi2019mediapipe, bazarevsky2020blazepose, topham2023hpesurvey}.
Table~\ref{tab:modalities} summarises the comparison that motivates the
markerless choice in this thesis.

\begin{table}[htbp]
\centering
\caption{Comparison of motion-analysis modalities for running-technique
measurement.}
\label{tab:modalities}
\renewcommand{\arraystretch}{1.2}
\begin{tabularx}{\textwidth}{l X X X}
\toprule
Property & Marker-based & Wearable sensors & Markerless video \\
\midrule
Instrumentation & Markers $+$ multi-camera rig & Body-worn IMUs & None (RGB video only) \\
Environment & Controlled lab & Field-capable & Any, incl.\ in-the-wild \\
Kinematic accuracy & Reference standard & High, drift-prone & Lower, noisier \\
Setup / calibration & High & Moderate & None \\
Applicable to internet video & No & No & Yes \\
Uncertainty signal & Physical calibration & Sensor model & Per-landmark confidence \\
\bottomrule
\end{tabularx}
\end{table}

The implication is direct: only markerless video can be applied retrospectively
to uncontrolled internet footage, so an in-the-wild system must accept the
attendant measurement noise and design around it rather than assume
laboratory-grade input.

% ---------------------------------------------------------------------
\section{Human Pose Estimation for Movement Understanding}
\label{sec:lit-pose}

Modern markerless pipelines rely on human pose estimators that localise a fixed
set of anatomical \emph{landmarks} (joints) per frame. For single-person,
monocular video, lightweight estimators produce \emph{2D landmarks} --- image
coordinates for each joint --- typically accompanied by a \emph{visibility} or
\emph{confidence} value that reflects the model's certainty that the joint is
present and correctly located
\cite{cao2019openpose, lugaresi2019mediapipe, bazarevsky2020blazepose}. This
per-landmark confidence is central to any downstream use: it allows low-quality
detections to be filtered or down-weighted, and it provides a measurable proxy
for input quality that can later be related to classification performance.

For running-technique analysis the \emph{lower-body} landmarks --- hips, knees,
ankles, heels, and toes --- are the most informative, because foot strike, knee
behaviour at touchdown, and ankle motion are precisely the descriptors that
distinguish good from poor form (Section~\ref{sec:lit-running})
\cite{novacheck1998biomechanics}. Unfortunately, these same landmarks are also
the least \emph{stable}. The feet and ankles are small, fast-moving, frequently
occluded by the opposite leg or by motion blur, and often near the image
boundary; their estimated positions therefore carry more noise and lower
confidence than torso landmarks
\cite{topham2023hpesurvey, panconi2024markerless}. This tension --- the most
diagnostic joints being the least reliably estimated --- is a defining challenge
of markerless running analysis and recurs throughout this thesis, both in the
choice of motion signals (Section~\ref{sec:lit-cycle}) and in the analysis of
error sources.

% ---------------------------------------------------------------------
\section{Running Cycle and Gait Phase Detection}
\label{sec:lit-cycle}

Because running is periodic, a natural pre-processing step is to segment the
continuous pose sequence into individual running cycles, conventionally bounded
by successive \emph{foot-strike} events of the same foot
(foot-strike-to-foot-strike) \cite{inostroza2025runningcycle, ghoussayni2004gait}.
Reliable cycle boundaries are a prerequisite for any phase-aligned
representation: if cycles are mis-bounded, samples drawn from "the same phase"
across different cycles will in fact correspond to different moments of the gait,
undermining comparison.

Two families of method are used to find these boundaries. \emph{Event-based}
methods detect foot-strike directly from a chosen motion signal, for instance by
locating peaks or valleys in a vertical-displacement or foot-contact proxy
signal \cite{ghoussayni2004gait, zhang2025heelstrike}. \emph{Periodicity-based}
methods first estimate the dominant period of the motion --- classically via
\emph{autocorrelation} of a one-dimensional signal --- and then use that period
to constrain or validate candidate boundaries, which makes detection more robust
to isolated spurious peaks \cite{cutting1978gaitperiodicity}. A practical
pipeline often combines the two: an autocorrelation estimate provides the
expected cycle length, and peak/valley detection on a foot-contact signal
supplies the actual boundaries, with the period acting as a consistency check.

The choice of signal interacts strongly with stability
(Section~\ref{sec:lit-pose}). Signals derived from torso or hip vertical motion
are smooth and easy to detect but only indirectly related to foot strike, while
signals derived from the noisy foot and heel landmarks are more directly tied to
contact events but more prone to error \cite{zhang2025heelstrike}. Selecting a
signal is therefore a trade-off between detection coverage and boundary
precision, a trade-off this thesis examines explicitly.

% ---------------------------------------------------------------------
\section{Pose-Based Action Classification}
\label{sec:lit-pose-class}

Once cycles are segmented and sampled, the resulting landmark sequences can be
classified by a range of neural architectures, ordered roughly by how much
temporal structure they model \cite{yan2018stgcn, jangua2021gait2dposes}. A
\emph{multilayer perceptron} (MLP) treats the flattened sequence as a
fixed-length vector and learns a static mapping; it is fast and low-variance but
ignores temporal ordering, which can be adequate when inputs are short and
well-aligned. A \emph{one-dimensional convolutional neural network} (1D-CNN)
applies temporal convolutions across the sequence, capturing local motion
patterns with relatively few parameters \cite{ordonez2016deephar}.

Recurrent models capture longer-range temporal dependencies. A
\emph{bidirectional long short-term memory} network (BiLSTM) processes the
sequence in both directions, allowing each phase to be interpreted in the
context of the whole cycle \cite{hochreiter1997lstm, schuster1997bidirectional}.
A \emph{CNN-BiLSTM} hybrid first extracts local features with convolutions and
then models their temporal evolution with a BiLSTM, combining local pattern
extraction with sequence modelling \cite{donahue2015lrcn}. More expressive
recurrent and hybrid models can represent richer dynamics, but on small or noisy
datasets their additional capacity also increases the risk of overfitting --- a
consideration that is especially relevant when the input cycles are themselves
imperfectly segmented.

% ---------------------------------------------------------------------
\section{Image-Based Video Classification}
\label{sec:lit-image}

An alternative to landmark sequences is to classify the \emph{appearance} of the
runner directly. Here each sampled phase is represented as an RGB crop centred
on the runner, and a convolutional backbone extracts visual features that a
temporal model then aggregates over the cycle. A widely used efficient design
pairs a lightweight pretrained convolutional backbone such as
\emph{MobileNetV2} with a recurrent temporal head such as an LSTM
(a MobileNetV2--LSTM model): the backbone encodes each frame into a feature
vector, and the LSTM models the sequence of per-frame features to produce a
cycle-level prediction \cite{sandler2018mobilenetv2, donahue2015lrcn}.

The appeal of image-based representations is that RGB crops retain information
absent from sparse landmarks --- body silhouette, limb contour, clothing-relative
motion, and surrounding visual context --- which could in principle aid
classification \cite{topham2023hpesurvey}. The cost is a substantially larger and
higher-variance input. Image pipelines must learn to ignore appearance variation
(background, clothing, lighting) that landmarks discard by construction, which
typically demands more training data; they are also sensitive to the stability
of the runner-centred crop, which itself depends on the same noisy pose estimates
discussed in Section~\ref{sec:lit-pose}.
Table~\ref{tab:representations} contrasts the two representation families.

\begin{table}[htbp]
\centering
\caption{Comparison of landmark-based and image-based cycle representations.}
\label{tab:representations}
\renewcommand{\arraystretch}{1.2}
\begin{tabularx}{\textwidth}{l X X}
\toprule
Property & Landmark-based & Image-based (RGB crop) \\
\midrule
Input per phase & Joint coordinates ($+$ confidence) & RGB crop (e.g.\ $224\times224$) \\
Dimensionality & Low & High \\
Information retained & Joint geometry only & Silhouette, contour, context \\
Invariance & Appearance-invariant by design & Must be learned \\
Typical model & MLP / CNN / BiLSTM / CNN-BiLSTM & CNN backbone $+$ temporal head \\
Data requirement & Lower & Higher \\
Main vulnerability & Pose-estimation noise & Crop instability, overfitting \\
\bottomrule
\end{tabularx}
\end{table}

% ---------------------------------------------------------------------
\section{Evaluation Protocols for Multi-Cycle Video Data}
\label{sec:lit-eval}

A subtlety specific to cycle-based video analysis is the choice of evaluation
unit. Because one video can yield several cycles, predictions may be scored at
the \emph{cycle level} --- treating each detected cycle as an independent sample
--- or at the \emph{video level} --- aggregating a video's cycles into a single
prediction by, for example, majority vote, mean probability, or
confidence-weighted vote \cite{singh2018gaitsurvey}. The two views answer
different questions: cycle-level metrics characterise per-cycle discriminability,
whereas video-level metrics estimate the performance a user would experience,
since the practical unit of interest is the video.

Two methodological hazards accompany this choice. First, if cycles from the same
video are allowed to appear in both training and test sets, the evaluation leaks
video-specific information and inflates scores; splitting strictly \emph{by
video} is therefore essential. Second, comparisons across
pipelines are only fair when the evaluation unit and protocol are held fixed; a
model evaluated on whole-video sequences is not directly comparable to one
evaluated per cycle. Sound practice is to report the split protocol explicitly,
to fix the evaluation unit within any comparison, and to state whether a tuned
decision threshold was used and how it was selected.

% ---------------------------------------------------------------------
\section{Research Gap}
\label{sec:lit-gap}

The reviewed literature establishes that running technique is a periodic,
lower-body-dominated phenomenon (Section~\ref{sec:lit-running}); that only
markerless video can be applied to in-the-wild footage, at the price of noisy
estimates (Section~\ref{sec:lit-modalities}); that the most diagnostic landmarks
are also the least stable (Section~\ref{sec:lit-pose}); and that cycle detection,
representation, and classification each offer well-studied options
(Sections~\ref{sec:lit-cycle}--\ref{sec:lit-image}). What remains
under-addressed is their combination on \emph{uncontrolled internet video} under
a rigorous, leakage-free evaluation.

Specifically, three gaps motivate this thesis. First, prior work seldom
\emph{separates} the contribution of automatic cycle detection from that of the
classifier; without a controlled upper bound, it is unclear whether failures
stem from segmentation or from learning. Second, landmark-based and image-based
representations are rarely compared on an \emph{identical} automatic cycle source
under a fixed split, so apparent representation advantages may merely reflect
protocol differences --- and historical, differently-evaluated milestones risk
being mistaken for controlled benchmarks. Third, the interaction between
pose/cycle quality and classification accuracy on in-the-wild data is rarely
quantified, leaving the dominant error source unidentified. This thesis addresses
these gaps with a matched handmade-upper-bound versus automatic-baseline design,
a controlled landmark-versus-image comparison on a shared cycle source, and an
explicit analysis of cycle-level versus video-level evaluation, threshold
selection, and pose/cycle quality
(Chapters~\ref{ch:methodology}--\ref{ch:discussion}).

% End of Chapter 2